\documentclass[11pt]{article}
\usepackage{amsmath,amssymb,amsthm,hyperref}
\newtheorem{theorem}{Theorem}
\title{Compact pole windows require a singular dictionary limit}
\date{24 November 2026}
\begin{document}
\maketitle

\section{A local obstruction}

Let $X$ be a nonempty finite query set and let $A(\theta)$ be a continuous
family of real $X\times H$ evaluation matrices. Its columns include every
freely weighted feature, so fixed polynomial correction features are charged
as additional columns. The coefficients multiplying these columns may be
arbitrarily large and depend arbitrarily on the desired labels.

\begin{theorem}
Suppose $G_0=A(\theta_0)^\top A(\theta_0)$ has positive determinant. Suppose
there is $z\in\mathbb R^X$ with every coordinate nonzero and
$A(\theta_0)^\top z=0$. Then there is an open neighborhood of $\theta_0$
on which no score $A(\theta)a$ has strict coordinate signs equal to those
of $z$.
\end{theorem}

\begin{proof}
Put $G(\theta)=A(\theta)^\top A(\theta)$ and define the polynomial expression
\[
 q(\theta)=\det G(\theta)z-
 A(\theta)\operatorname{adj}(G(\theta))A(\theta)^\top z.
\]
The adjugate identity gives $A(\theta)^\top q(\theta)=0$ for every
$\theta$, including singular matrices. At the center,
$q(\theta_0)=\det G_0 z$. By continuity all coordinates of $q(\theta)$
have the corresponding signs of $z$ throughout some common neighborhood.
If $A(\theta)a$ had these strict signs, then
$q(\theta)^\top A(\theta)a$ would be strictly positive, contradicting
its being zero. No bound on $a$ was used.
\end{proof}

This theorem is kernel checked as
\texttt{HeadComplexity.AxisZeroGluing.jointFamily\_compact\_regular\_dictionary\_obstruction}
in
\texttt{formalization/HeadComplexity/Results/JointFamilyCompactRank.lean}.
The auxiliary strict dual lemma and adjugate annihilation identity are
also checked. The hypotheses are explicit; the file does not infer genericity
or construct a nonzero dual vector from a rank count.

\section{Consequences for the proposed local decoder}

For fixed positive row directions, a reciprocal column has the form
\[
 A_{x,h}(p)=\frac{1}{p_h+\sum_i w_{h,i}(1-x_i)},\qquad w_{h,i}>0.
\]
Positive poles give positive denominators on the entire row cube. For
independently chosen column poles in the reverse cube-star construction,
positivity on multihot column inputs generally forces the permitted pole
variation to shrink with the number of columns. Merely reusing a fixed
regular dictionary with a narrower window cannot eventually yield all
strict label tables: the theorem excludes one fixed row sign pattern
throughout a sufficiently small window. Allowing arbitrary residues,
permuting the poles, or adding a fixed finite correction space does not
remove this obstruction when the enlarged center matrix still satisfies
the displayed hypotheses.

If the center span contains no coordinate vector and is proper, elementary
linear algebra yields a full-support normal $z$: for each coordinate the
normal space is not contained in its zero-coordinate hyperplane, and a
real vector space is not the union of finitely many proper subspaces.
This consequence is a written proof, not an additional Lean declaration.

\section{Scope and surviving positive routes}

A compact decoder factory is allowed to choose its row directions as a
function of the permitted window. The theorem therefore does not rule out
families deliberately approaching a singular center. In particular, when
all directions approach the same vector, the limiting evaluation matrix
can have rank at most the number of Hamming weights, leaving a large
kernel if $H>k+1$. Pole/residue permutations can then select different
higher-order jets after common low-order corrections. Proving that these
residual branches meet every required strict orthant remains open.

Nor does the theorem exclude a finite window of fixed width, a dictionary
with a singular center, or nonlinear correction features whose own
parameters vary with the pole window. It supplies no lower bound for the
original unrestricted worst-case head complexity.

\end{document}
